\subsection{A distribution-independent reserve transformation}\label{sec:reserve-general}
Let $M_*$ be any Borel, jointly feasible, pointwise DSIC and ex-post IR
mechanism on $[0,1]^4$, with allocations in $[0,1]^2$ and bounded expected
payments. Assume that its selected allocation excludes zero-valued items,
\[
 z_{ij}=0\quad\Longrightarrow\quad x^*_{ij}(z)=0
 \qquad\text{at every profile, including ties}.
\]
Borel measurability and bounded payments ensure that clamping onto boundary
faces preserves measurability and integrability. The incentive argument
itself uses no distributional assumption and no finite allocation range.
For $0\le\tau<1$ put $s_\tau=1-\tau$ and
$T_\tau(t)=\max\{(t-\tau)/s_\tau,0\}$, coordinatewise on the full profile.
Define
\begin{equation}\label{eq:reserve-transform}
 x^\tau_{ij}(v)=x^*_{ij}(T_\tau v),\qquad
 p_i^\tau(v)=s_\tau p_i^*(T_\tau v)+\tau\sum_jx^*_{ij}(T_\tau v).
\end{equation}
Use the seed's given allocation and payment selection at the transformed
profile, including its tie rule. There is no additional allocation gate.

\begin{proposition}[Reserve preservation]\label{prop:reserve-preservation}
For every $\tau\in[0,1)$, \eqref{eq:reserve-transform} defines a Borel,
pointwise DSIC, ex-post IR and jointly feasible mechanism on $[0,1]^4$.

\end{proposition}
\begin{proof}
Fix an arbitrary opponent report, a true type $v_i$, and a report $r_i$.
Its deviation utility is bounded by
\begin{align*}
 &\sum_j(v_{ij}-\tau)x^*_{ij}(T_\tau r_i,T_\tau v_{-i})
       -s_\tau p_i^*(T_\tau r_i,T_\tau v_{-i})\\
 &\quad\le s_\tau\big[T_\tau v_i\cdot x_i^*(T_\tau r_i,T_\tau v_{-i})
       -p_i^*(T_\tau r_i,T_\tau v_{-i})\big]\\
 &\quad\le s_\tau u_i^*(T_\tau v).
\end{align*}
The first inequality uses nonnegative allocations; the second is DSIC
of $M_*$. At truth the first inequality is equality by
the zero-value-exclusion hypothesis, including $v_{ij}=\tau$ and $v_{ij}<\tau$.
Hence truthful utility is $s_\tau u_i^*(T_\tau v)\ge0$.
Both allocations come from the same feasible transformed profile, so
capacity holds pointwise. The map is continuous and the selected base
rule is Borel. Bounded seed payments keep transformed payments bounded,
so integrability also holds. All assertions include exceptional reports and ties and
use the expected-utility convention of Section~\ref{sec:model}.
\end{proof}


\begin{figure}[t]
\centering\small
\fbox{\begin{minipage}{0.92\linewidth}
\centering
Reports $v\in[0,1]^4$
$\ \xrightarrow{\ T_\tau\ }\ $
one transformed profile $z$
$\ \xrightarrow{\ M_*\ }\ $
$(x^*(z),p^*(z))$
\[
 x^\tau(v)=x^*(z),\qquad
 p_i^\tau(v)=(1-\tau)p_i^*(z)+\tau\sum_jx^*_{ij}(z).
\]
Both bidders use the same feasible allocation. Zero-valued transformed
items receive zero allocation. Under independent uniform reports,
each coordinate has law $\tau\delta_0+(1-\tau)\operatorname{Unif}[0,1]$.
\end{minipage}}
\caption{One common report transformation preserves supply and creates the
boundary atoms used in the exact revenue calculation. The seed is fully
specified by the three menus in Section~\ref{sec:refined-mechanism}.}
\label{fig:reserve-map}
\end{figure}
